\subsection{Matrices and Graph Theory}

\begin{remarkbox}[title={\bfseries Enrichment Material}]
The topics in this appendix go \textbf{beyond} the Queensland (QCAA) Specialist Mathematics syllabus. They are included to give curious students a preview of the linear algebra they will meet in first- and second-year university courses in mathematics, engineering, computer science, data science, and AI. None of this material is examinable in QLD Year~12 assessments.
\end{remarkbox}

This section connects two areas of mathematics: \textbf{matrices} and \textbf{graph theory}. A graph represents relationships between objects, while a matrix provides a compact numerical way to store and analyse those relationships.

This connection is widely used in computer science, transport networks, social networks and web search.

\subsubsection*{Graphs}

A \textbf{graph} consists of:
\begin{itemize}[leftmargin=*]
    \item \textbf{vertices} (or nodes): the objects;
    \item \textbf{edges}: connections between pairs of vertices.
\end{itemize}

For example, suppose four railway stations $A, B, C, D$ are connected as follows:
\[ A-B, \qquad A-C, \qquad B-C, \qquad C-D. \]
We can think of the vertices as stations and the edges as direct railway connections.

\subsubsection*{The Adjacency Matrix}

A graph can be represented by an \textbf{adjacency matrix}.

Let
\[
M_{ij} =
\begin{cases}
1, & \text{if vertices } i \text{ and } j \text{ are connected}, \\
0, & \text{otherwise}.
\end{cases}
\]

For our railway network,
\[
M = 
\begin{bmatrix}
0 & 1 & 1 & 0 \\
1 & 0 & 1 & 0 \\
1 & 1 & 0 & 1 \\
0 & 0 & 1 & 0
\end{bmatrix}.
\]

The rows and columns correspond to $A, B, C, D$. 

For example,
\[ M_{AC} = 1 \]
because $A$ and $C$ are directly connected, while
\[ M_{AD} = 0 \]
because there is no direct connection between $A$ and $D$.

Because this graph is \textbf{undirected}, its adjacency matrix is symmetric:
\[ M = M^T. \]

\subsubsection*{Matrix Powers Reveal Routes}

Here is the surprising connection.

If $M$ is an adjacency matrix, then the entries of
\[ M^2 = M \times M \]
count the number of \textbf{two-edge walks} between vertices.

For example, looking at the entry for $A$ to $D$:
\[ (M^2)_{AD} = M_{AA}M_{AD} + M_{AB}M_{BD} + M_{AC}M_{CD} + M_{AD}M_{DD} \]
Therefore,
\[ (M^2)_{AD} = (0)(0) + (1)(0) + (1)(1) + (0)(0) = 1. \]

So there is exactly one two-edge route from $A$ to $D$:
\[ A \rightarrow C \rightarrow D. \]

More generally,
\[ \boxed{ (M^k)_{ij} = \text{number of walks of length } k \text{ from } i \text{ to } j. } \]

This gives matrix multiplication a concrete interpretation: multiplying matrices can reveal connectivity in a network. In fact, this example directly reinforces why matrix multiplication is defined as a row-by-column dot product!

\begin{remark}[Note]
A \emph{walk} may revisit vertices or edges, so it is not necessarily a simple path.
\end{remark}

\subsubsection*{Why It Matters}

This matrix--graph connection appears in many real systems:
\begin{itemize}[leftmargin=*]
    \item \textbf{Transport:} stations are vertices and routes are edges.
    \item \textbf{Social networks:} people are vertices and friendships or follows are edges.
    \item \textbf{Computer networks:} devices are vertices and communication links are edges.
    \item \textbf{The Web:} webpages are vertices and hyperlinks form directed edges.
\end{itemize}

Large graphs can contain millions or billions of connections. Representing them using matrices allows computers to apply linear algebra to analyse these networks.

\vspace{1em}
\begin{center}
\fbox{\textbf{Graph} $\longleftrightarrow$ \textbf{Adjacency Matrix} $\longleftrightarrow$ \textbf{Matrix Algebra}}
\end{center}
\vspace{0.5em}

\textbf{Takeaway:} Graph theory describes \textbf{who is connected to whom}; matrices turn those connections into numbers that we can calculate with.
